\documentclass[12pt]{article}
\usepackage[utf8]{inputenc}
\usepackage[T1]{fontenc}
\usepackage[margin=1in]{geometry}
\usepackage{setspace}
\usepackage{xcolor}
\usepackage[colorlinks=true, urlcolor=blue!60!black, linkcolor=blue!60!black]{hyperref}
\setlength{\parindent}{2em}
\setlength{\parskip}{1.25em}
\renewcommand{\baselinestretch}{1.0}
\begin{document}
\begin{center}
{\LARGE \textit{Champion v. Ames}, 188 U.S. 321 (1903)\par}
\vspace{0.8em}
{\large Case Brief --- Four-Page Edition\par}
\end{center}
\vspace{1em}

\section{Memory Jogger}

In the Lottery Case, the Court upheld 5--4 the federal ban on carrying lottery tickets across state lines, holding that lottery tickets are articles of commerce and that the power to regulate interstate commerce includes the power to prohibit it, the root of the federal ``police power'' exercised through the channels of commerce (\href{https://supreme.justia.com/cases/federal/us/188/321/}{Justia}).

\section{Detailed Case Facts}

After state prohibitions failed to stop out-of-state lottery operators using national transportation networks, and after barring lottery materials from the mails and from importation, Congress enacted the Federal Lottery Act of 1895, ``An Act for the Suppression of Lottery Traffic through National and Interstate Commerce and the Postal Service.'' Section~1 made it a crime to cause any lottery ticket to be ``carried from one state to another in the United States,'' punishable by up to two years' imprisonment and a \$1,000 fine.

In 1899, Charles F. Champion and two others were indicted in the Northern District of Texas for conspiring to violate the Act. They had deposited a box containing two tickets of the Pan-American Lottery Company, which purported to hold monthly drawings in Paraguay for a \$32,000 prize, with the Wells-Fargo Express Company in Dallas for shipment by rail to Fresno, California.

\section{Procedural History}

Arrested in Chicago on a federal removal warrant, Champion was committed to the custody of U.S. Marshal John C. Ames and sought habeas corpus in the Circuit Court for the Northern District of Illinois, arguing that the Act exceeded the commerce power. The court dismissed the writ, and Champion appealed directly. The case was argued three times before Justice Holmes joined the Court, and it was decided on February 23, 1903.

\section{Judicial Votes}

Affirmed, \textbf{5--4}.
\begin{itemize}
\item \textbf{Majority:} Justice Harlan, joined by Justices Brown, White, McKenna, and Holmes.
\item \textbf{Concurrences:} None.
\item \textbf{Dissent:} Chief Justice Fuller, joined by Justices Brewer, Shiras, and Peckham.
\end{itemize}

\section{Holding}

Lottery tickets are subjects of traffic and therefore of commerce, and their carriage between States by independent carriers is interstate commerce. Congress's plenary power to regulate interstate commerce includes the power to prohibit the interstate transportation of articles it deems injurious to public morals, and doing so violates neither the Tenth Amendment nor the Fifth Amendment's Due Process Clause.

\section{Analysis of Opinions}

\textbf{Majority (Harlan): commerce.} Relying on \textit{Gibbons v. Ogden}, Harlan held that commerce includes ``navigation, intercourse, communication, traffic, the transit of persons, and the transmission of messages by telegraph.'' The tickets represented a right to a prize, had a market price, and were carried for hire, so they were articles of commerce.

\textbf{Regulation includes prohibition.} Under \textit{McCulloch}, Congress has wide discretion over means; it need not merely regulate the manner of shipping an enterprise the Court had condemned as a ``widespread pestilence.'' Congress already banned interstate shipment of diseased cattle, prohibited restraints of trade under the Sherman Act, and excluded liquor from dry States under the Wilson Act.

\textbf{Tenth and Fifth Amendments.} Exercising a delegated power cannot invade reserved powers, and Congress had only supplemented state anti-lottery policies without reaching purely internal sales. No one has a liberty-of-contract right to inject into commerce an article ``confessedly injurious to the public morals.'' Harlan declined to map the power's limits in advance; for unwise legislation, the check is the political process.

\textbf{Fuller (dissenting).} The Act's avowed purpose was to suppress lotteries, a local matter for the States' police power under \textit{E. C. Knight}; Congress could exclude lottery matter from the mails but not from private carriage, and upholding the law created a national police power. Like insurance contracts in \textit{Paul v. Virginia}, tickets are not articles of commerce. The commerce power was meant to free interstate trade from state discrimination, ``not to transfer the power of restriction''; diseased cattle infect other goods, but lottery tickets ``cannot communicate bad principles by contact.''

\textbf{Later developments.} The principle sustained the Pure Food and Drug Act, the Mann Act, and the Motor Vehicle Theft Act. \textit{Hammer v. Dagenhart} (1918) tried to limit it to goods harmful in themselves, but \textit{United States v. Darby} (1941) overruled \textit{Hammer}. \textit{Champion} underlies the ``channels'' category of \textit{United States v. Lopez}.

\section{Examples: Future Applications}

\begin{enumerate}
\item \textbf{Same-Side (Upheld): Interstate Medication Ban.} Congress bans interstate shipment of an abortifacient lawful in both States. Under \textit{Champion} and \textit{Darby}, the channels power reaches the medication regardless of motive.
\item \textbf{Same-Side (Upheld): Unserialized Gun Kits.} A ban on shipping build-it-yourself firearm kits across state lines regulates articles of traffic; crime-control motive is irrelevant.
\item \textbf{Opposite-Side (External Limits): The Brochures.} The same statute's ban on shipping instructional brochures implicates the First Amendment, and Congress here overrides both States' policies rather than supplementing them.
\item \textbf{Opposite-Side (No Channel): Public Intoxication.} A federal ban on public intoxication with no interstate element regulates no channel and fails under \textit{Lopez} and \textit{Morrison}.
\item \textbf{Fence-Sitter: Possession After Interstate Travel.} A ban on possessing any vaping device that once crossed state lines invokes \textit{Scarborough}'s minimal nexus, but \textit{Champion} rested on stopping articles in transit.
\end{enumerate}

\section{Critique}

\textbf{The ``pollution'' metaphor and scope creep.} Harlan leaned on the moral character of lotteries, suggesting commerce could be protected from ``pollution.'' As Fuller foresaw, once Congress may close channels on moral or policy grounds, the ``harmful-in-itself'' limit collapses, producing the whiplash from \textit{Champion} to \textit{Hammer} to \textit{Darby}.

\textbf{Formalism versus pretext.} Harlan's claim that the Act left ``completely internal'' affairs untouched is formally accurate but functionally illusory: because national distribution is the lifeblood of most enterprises, banning interstate shipment lets Congress regulate production and local behavior without claiming a police power.

\section{Key Quotations}

\begin{quote}
``As a state may, for the purpose of guarding the morals of its own people, forbid all sales of lottery tickets within its limits, so Congress, for the purpose of guarding the people of the United States against the `widespread pestilence of lotteries' and to protect the commerce which concerns all the states, may prohibit the carrying of lottery tickets from one state to another.''
\end{quote}
\noindent --- Justice Harlan, for the Court, 188 U.S. at 357.

\begin{quote}
``We are of opinion that lottery tickets are subjects of traffic, and therefore are subjects of commerce, and the regulation of the carriage of such tickets from state to state, at least by independent carriers, is a regulation of commerce among the several states.''
\end{quote}
\noindent --- Justice Harlan, for the Court, 188 U.S. at 354.

\begin{quote}
``If a state, when considering legislation for the suppression of lotteries within its own limits, may properly take into view the evils that inhere in the raising of money in that mode, why may not Congress, invested with the power to regulate commerce among the several states, provide that such commerce shall not be polluted by the carrying of lottery tickets from one state to another?''
\end{quote}
\noindent --- Justice Harlan, for the Court, 188 U.S. at 356.

\end{document}
